\addtocontents{toc}{\protect\setcounter{tocdepth}{0}} % Disable section numbering in ToC for appendices

\chapter{LAMPIRAN}\label{ch:appendix}
\vspace*{20pt}

% ========================= START ============================
% Appendix Section Formatting, only move/modify cautiously.
% ============================================================
% Custom counter for appendix sections
\newcounter{appendixsection}
\setcounter{appendixsection}{0}

% Redefine \thesection to use our custom counter
\renewcommand{\thesection}{\arabic{appendixsection}}

% Modify \section to increment counter and set label properly
\makeatletter
\let\oldsection\section
\renewcommand{\section}{%
  \@ifnextchar[%]
    {\appendix@section@opt}
    {\appendix@section@noopt}%
}

\def\appendix@section@opt[#1]#2{%
  \refstepcounter{appendixsection}%
  \oldsection[#1]{#2}%
}

\def\appendix@section@noopt#1{%
  \refstepcounter{appendixsection}%
  \oldsection{#1}%
}
\makeatother

% Custom formatting for section titles
\titleformat{\section}[block]{\bfseries\normalsize}{}{0pt}{Lampiran \thesection: }
\titlespacing*{\section}{0pt}{0.5em}{0.5em}

% Ensure each section (except the first) starts on a new page
\newcounter{sectioncount}
\setcounter{sectioncount}{0}

\let\oldsectiontwo\section
\renewcommand{\section}[1]{%
  \ifnum\thesectioncount>0
    \clearpage
  \fi
  \stepcounter{sectioncount}%
  \oldsectiontwo{#1}%
}

% Disable paragraph indentation in appendix
\setlength{\parindent}{0pt}

% ============================================================
% Appendix Section Formatting, the following are appendix contents.
% =========================== END ============================

\section{Example And Rules in Appendix}\label{apx:example_and_rules}
\mylampiran{\textit{Example and rules in Appendix}}

\begin{comment}
You can use this type of comment to provide additional information or instructions for the appendix.
Will not be displayed in the final document.
\end{comment}

\begin{longtable}{|>{\arraybackslash}p{1cm}|>{\centering\arraybackslash}p{2.5cm}|>{\centering\arraybackslash}p{2.5cm}|}

    % Title Header
    \hline
    \multicolumn{3}{|c|}{\textbf{A sample table}} \\
    \hline

    % Subgroup 1
    \textbf{No} & \multicolumn{2}{|l|}{\textbf{Sub-group 1}} \\
    \hline
    & $Column_{1}$ & $Column_{2}$\\
    \hline
    \endfirsthead

    % Header repeated on continuation pages
    % \hline
    % \multicolumn{5}{|r|}{\textit{Continue\ldots}} \\
    \hline
    & $Column_{1}$ & $Column_{2}$\\
    \hline
    \endhead

    \hline
    \endfoot

    \csvreader[
        no head,
        late after line=\\\hline
    ]{Back_matters/attachments/apx-sample-1.csv}{}
    {\csvcoli & \csvcoliii}

    % Subgroup 2
    \textbf{Lingkungan:} & \multicolumn{2}{|l|}{\textbf{Sub-group 2}} \\
    \hline
    \csvreader[
        no head,
        late after line=\\\hline
    ]{Back_matters/attachments/apx-sample-2.csv}{}
    {\csvcoli & \csvcoliii}

\end{longtable}

Pelatihan dilakukan dari episode 1 hingga episode 750 dengan pencatatan kesimpulan metrik pada setiap akhir episode.
Terdapat banyak metrik yang dicatat, namun tabel performa ini hanya menyajikan fokus metrik utama yang dibahas.
Di mana Group adalah konfigurasi model terbaik dalam kelompoknya.
Lalu $CCR_{e=E}$ yang merupakan \acrlong{ccr} (\acrshort{ccr}) pada episode terakhir (episode ke-750).
Kemudian, $CCR_{avg}$ yaitu rata-rata \acrshort{ccr} selama 750 episode.
Ada juga $CCR_{max}$ yaitu nilai maksimum \acrshort{ccr} yang dicapai selama 750 episode.
Terakhir, $Y_{cum}$ adalah jumlah kumulatif \acrlong{y} (\acrshort{y}) yang dicapai selama 750 episode.
Tabel ini dikelompokkan lagi berdasarkan konfigurasi lingkungan dan jenis seed yang digunakan selama pelatihan awal.


\section{Source Code and Data}\label{apx:source_code_and_data}
Use this appendix to provide source code listings, data files, or any other supplementary material that is relevant to your document. 
You can include code snippets, scripts, or even entire files if necessary. 
Make sure to format the code properly for readability.